\documentclass[a4paper,11pt]{article}

\usepackage[a4paper,margin=2cm]{geometry}
\usepackage[T1]{fontenc}
\usepackage[utf8]{inputenc}
\usepackage[ngerman,english]{babel}
\usepackage{lmodern}
\usepackage{microtype}
\usepackage{xcolor}
\usepackage{parskip}
\usepackage{enumitem}
\usepackage[most]{tcolorbox}
\usepackage{titlesec}
\usepackage[hidelinks]{hyperref}

\definecolor{HeaderBlueDark}{RGB}{70,110,170}
\definecolor{RowBlueLight}{RGB}{235,243,255}
\definecolor{RowBlueDark}{RGB}{220,233,250}
\definecolor{SoftGray}{RGB}{90,90,90}

\setlength{\parindent}{0pt}
\setlist[itemize]{leftmargin=1.4em,itemsep=0.35em,topsep=0.35em}
\linespread{1.06}

\titleformat{\section}[block]
{\Large\bfseries\color{HeaderBlueDark}}
{}{0pt}{}
\titlespacing{\section}{0pt}{1.3em}{0.8em}

\titleformat{\subsection}[block]
{\large\bfseries\color{HeaderBlueDark}}
{}{0pt}{}
\titlespacing{\subsection}{0pt}{1.0em}{0.6em}

\newtcolorbox{exbox}{enhanced,breakable,colback=RowBlueLight,colframe=RowBlueDark,boxrule=0.4pt,arc=1.5mm,left=1.2mm,right=1.2mm,top=1mm,bottom=1mm,title={Beispiele \textnormal{(Examples)}},fonttitle=\bfseries,colbacktitle=RowBlueDark,coltitle=black}

\NewDocumentEnvironment{grammarentry}{mm+b}{%
    \begin{tcolorbox}[enhanced,breakable,colback=white,colframe=HeaderBlueDark,boxrule=0.6pt,arc=1.5mm,left=1.5mm,right=1.5mm,top=1mm,bottom=1mm,colbacktitle=HeaderBlueDark,coltitle=white,fonttitle=\bfseries,title={#1\hfill\normalfont\footnotesize #2}]
    #3
    \end{tcolorbox}
}{}

\newcommand{\GermanText}[1]{\textbf{Deutsch: }#1\par}
\newcommand{\EnglishText}[1]{{\small\color{SoftGray}\textbf{English: }#1\par}}
\newcommand{\ExampleItem}[2]{\item \textbf{#1}\\{\small\color{SoftGray}#2}}

\title{Platzhalter und Korrelat im Deutschen}
\date{}

\begin{document}
    \maketitle
    \tableofcontents
    \newpage
    
    
    \section{Grundidee}
        
        \begin{grammarentry}{Platzhalter / Korrelat}{Grundbegriff}
            \GermanText{Ein Platzhalter oder Korrelat ist ein Wort, das im Hauptsatz steht und auf einen späteren Nebensatz oder Infinitivsatz vorbereitet.}
            \EnglishText{A placeholder or correlative is a word in the main clause that prepares for a following subordinate clause or infinitive clause.}
            
            Häufige Korrelate sind:
            \begin{itemize}
                \item \textbf{es}
                \item \textbf{darauf}
                \item \textbf{daran}
                \item \textbf{damit}
                \item \textbf{dafür}
                \item \textbf{dazu}
                \item \textbf{darüber}
                \item \textbf{darum}
            \end{itemize}
            
            Wichtig: Das Korrelat verweist oft nicht auf ein einzelnes Nomen, sondern auf eine ganze Handlung, einen ganzen Infinitivsatz oder einen ganzen Nebensatz.
        \end{grammarentry}
        
        \begin{exbox}
            \begin{itemize}
                \ExampleItem{Ich bevorzuge es, früh aufzustehen.}{I prefer to get up early.}
                \ExampleItem{Ich freue mich darauf, dich zu sehen.}{I look forward to seeing you.}
                \ExampleItem{Es ist wichtig, Deutsch zu lernen.}{It is important to learn German.}
            \end{itemize}
        \end{exbox}
    
    
    \section{Das Korrelat \textit{es}}
        
        \begin{grammarentry}{es als Platzhalter}{bei Infinitivsatz und dass-Satz}
            \GermanText{\textbf{Es} kann auf einen späteren Infinitivsatz oder dass-Satz verweisen.}
            \EnglishText{\textbf{Es} can refer forward to a later infinitive clause or dass-clause.}
            
            Struktur:
            \begin{itemize}
                \item \textbf{Ich finde es gut, ... zu ...}
                \item \textbf{Ich halte es für wichtig, ... zu ...}
                \item \textbf{Ich bevorzuge es, ... zu ...}
                \item \textbf{Es ist wichtig, dass ...}
                \item \textbf{Es freut mich, dass ...}
            \end{itemize}
        \end{grammarentry}
        
        \begin{exbox}
            \begin{itemize}
                \ExampleItem{Ich finde es gut, jeden Tag Deutsch zu üben.}{I think it is good to practice German every day.}
                \ExampleItem{Ich halte es für wichtig, pünktlich zu sein.}{I consider it important to be on time.}
                \ExampleItem{Ich bevorzuge es, in einem ruhigen Land zu leben.}{I prefer to live in a quiet country.}
                \ExampleItem{Es ist schade, dass du nicht kommen kannst.}{It is a pity that you cannot come.}
            \end{itemize}
        \end{exbox}
    
    
    \section{\textit{es} und der Infinitivsatz}
        
        \begin{grammarentry}{Ich bevorzuge es, ... zu ...}{wichtige Struktur}
            \GermanText{Bei einigen Verben benutzt man oft \textbf{es}, wenn danach ein Infinitivsatz mit \textbf{zu} kommt.}
            \EnglishText{With some verbs, \textbf{es} is often used when a following infinitive clause with \textbf{zu} comes after it.}
            
            Beispiel:
            \begin{itemize}
                \item \textbf{Ich bevorzuge es, allein zu arbeiten.}
                \item Hier steht \textbf{es} für: \textbf{allein zu arbeiten}
            \end{itemize}
            
            Das heißt:
            \begin{itemize}
                \item \textbf{Ich bevorzuge es, allein zu arbeiten.}
                \item = \textbf{Ich bevorzuge allein zu arbeiten.}
            \end{itemize}
            
            Die Version mit \textbf{es} klingt oft natürlicher und klarer.
        \end{grammarentry}
        
        \begin{exbox}
            \begin{itemize}
                \ExampleItem{Ich bevorzuge es, morgens Kaffee zu trinken.}{I prefer to drink coffee in the morning.}
                \ExampleItem{Sie liebt es, neue Sprachen zu lernen.}{She loves learning new languages.}
                \ExampleItem{Er hasst es, lange zu warten.}{He hates waiting for a long time.}
            \end{itemize}
        \end{exbox}
    
    
    \section{Das Korrelat bei Adjektiven}
        
        \begin{grammarentry}{es ist + Adjektiv + zu / dass}{sehr häufig}
            \GermanText{Viele Sätze mit Adjektiven benutzen \textbf{es} als Platzhalter.}
            \EnglishText{Many sentences with adjectives use \textbf{es} as a placeholder.}
            
            Struktur:
            \begin{itemize}
                \item \textbf{Es ist wichtig, ... zu ...}
                \item \textbf{Es ist möglich, ... zu ...}
                \item \textbf{Es ist schwer, ... zu ...}
                \item \textbf{Es ist schön, dass ...}
                \item \textbf{Es ist klar, dass ...}
            \end{itemize}
        \end{grammarentry}
        
        \begin{exbox}
            \begin{itemize}
                \ExampleItem{Es ist wichtig, regelmäßig zu lernen.}{It is important to study regularly.}
                \ExampleItem{Es ist schwer, diese Regel sofort zu verstehen.}{It is difficult to understand this rule immediately.}
                \ExampleItem{Es ist klar, dass du mehr üben musst.}{It is clear that you have to practice more.}
            \end{itemize}
        \end{exbox}
    
    
    \section{\textit{es} kann manchmal wegfallen}
        
        \begin{grammarentry}{Infinitivsatz am Anfang}{ohne es}
            \GermanText{Wenn der Infinitivsatz am Anfang steht, fällt \textbf{es} meistens weg.}
            \EnglishText{When the infinitive clause is at the beginning, \textbf{es} usually disappears.}
            
            Vergleich:
            \begin{itemize}
                \item \textbf{Es ist wichtig, Deutsch zu lernen.}
                \item \textbf{Deutsch zu lernen ist wichtig.}
            \end{itemize}
        \end{grammarentry}
        
        \begin{exbox}
            \begin{itemize}
                \ExampleItem{Es ist gesund, viel Wasser zu trinken.}{It is healthy to drink a lot of water.}
                \ExampleItem{Viel Wasser zu trinken ist gesund.}{Drinking a lot of water is healthy.}
                \ExampleItem{Es macht Spaß, Deutsch zu sprechen.}{It is fun to speak German.}
                \ExampleItem{Deutsch zu sprechen macht Spaß.}{Speaking German is fun.}
            \end{itemize}
        \end{exbox}
    
    
    \section{Korrelate mit Präpositionen}
        
        \begin{grammarentry}{darauf, daran, damit, dafür, dazu}{Präposition + da}
            \GermanText{Wenn ein Verb oder Adjektiv eine feste Präposition braucht, benutzt man oft ein Korrelat mit \textbf{da-}.}
            \EnglishText{When a verb or adjective needs a fixed preposition, German often uses a correlative with \textbf{da-}.}
            
            Beispiele:
            \begin{itemize}
                \item \textbf{sich freuen auf} $\rightarrow$ \textbf{darauf}
                \item \textbf{denken an} $\rightarrow$ \textbf{daran}
                \item \textbf{rechnen mit} $\rightarrow$ \textbf{damit}
                \item \textbf{sorgen für} $\rightarrow$ \textbf{dafür}
                \item \textbf{beitragen zu} $\rightarrow$ \textbf{dazu}
                \item \textbf{sich ärgern über} $\rightarrow$ \textbf{darüber}
            \end{itemize}
        \end{grammarentry}
        
        \begin{exbox}
            \begin{itemize}
                \ExampleItem{Ich freue mich darauf, dich zu sehen.}{I look forward to seeing you.}
                \ExampleItem{Ich denke daran, morgen einzukaufen.}{I am thinking about going shopping tomorrow.}
                \ExampleItem{Ich rechne damit, dass er später kommt.}{I expect that he will come later.}
                \ExampleItem{Ich sorge dafür, dass alles fertig ist.}{I make sure that everything is ready.}
                \ExampleItem{Das trägt dazu bei, dass wir schneller lernen.}{That contributes to us learning faster.}
            \end{itemize}
        \end{exbox}
    
    
    \section{Warum heißt es \textit{darauf} und nicht nur \textit{auf}?}
        
        \begin{grammarentry}{da + Präposition}{Verweis auf einen ganzen Satz}
            \GermanText{\textbf{Darauf} bedeutet hier nicht immer wörtlich on it. Es verweist oft auf einen ganzen späteren Satz.}
            \EnglishText{\textbf{Darauf} does not always literally mean on it. It often refers to an entire following clause.}
            
            Beispiel:
            \begin{itemize}
                \item \textbf{Ich freue mich darauf, dich zu sehen.}
                \item \textbf{darauf} = \textbf{dich zu sehen}
            \end{itemize}
            
            Ohne \textbf{darauf} hat der Satz oft eine andere Bedeutung oder klingt unvollständig.
            
            Vergleich:
            \begin{itemize}
                \item \textbf{Ich freue mich, dich zu sehen.}
                \item Bedeutung: Ich bin jetzt froh, weil ich dich sehe.
                \item \textbf{Ich freue mich darauf, dich zu sehen.}
                \item Bedeutung: Ich freue mich auf etwas in der Zukunft.
            \end{itemize}
        \end{grammarentry}
    
    
    \section{Korrelat mit dass-Satz}
        
        \begin{grammarentry}{Hauptsatz + Korrelat + dass-Satz}{häufige Struktur}
            \GermanText{Ein Korrelat kann auch auf einen \textbf{dass-Satz} verweisen.}
            \EnglishText{A correlative can also refer to a \textbf{dass-clause}.}
            
            Struktur:
            \begin{itemize}
                \item \textbf{Ich finde es gut, dass ...}
                \item \textbf{Ich rechne damit, dass ...}
                \item \textbf{Ich achte darauf, dass ...}
                \item \textbf{Ich bin froh darüber, dass ...}
            \end{itemize}
        \end{grammarentry}
        
        \begin{exbox}
            \begin{itemize}
                \ExampleItem{Ich finde es gut, dass du jeden Tag übst.}{I think it is good that you practice every day.}
                \ExampleItem{Ich rechne damit, dass der Zug Verspätung hat.}{I expect that the train will be delayed.}
                \ExampleItem{Ich achte darauf, dass ich keine Fehler mache.}{I pay attention to not making mistakes.}
                \ExampleItem{Ich bin froh darüber, dass du hier bist.}{I am happy that you are here.}
            \end{itemize}
        \end{exbox}
    
    
    \section{Korrelat mit Infinitivsatz}
        
        \begin{grammarentry}{Hauptsatz + Korrelat + Infinitivsatz}{zu + Infinitiv}
            \GermanText{Ein Korrelat kann auf einen Infinitivsatz mit \textbf{zu} verweisen.}
            \EnglishText{A correlative can refer to an infinitive clause with \textbf{zu}.}
            
            Struktur:
            \begin{itemize}
                \item \textbf{Ich habe vor, ... zu ...}
                \item \textbf{Ich freue mich darauf, ... zu ...}
                \item \textbf{Ich achte darauf, ... zu ...}
                \item \textbf{Ich versuche es, ... zu ...}
            \end{itemize}
            
            Wichtig: Bei Infinitivsätzen steht \textbf{zu} direkt vor dem Infinitiv.
        \end{grammarentry}
        
        \begin{exbox}
            \begin{itemize}
                \ExampleItem{Ich freue mich darauf, nach Graz zu kommen.}{I look forward to coming to Graz.}
                \ExampleItem{Ich achte darauf, langsam zu sprechen.}{I pay attention to speaking slowly.}
                \ExampleItem{Ich finde es schwer, schnell Deutsch zu sprechen.}{I find it difficult to speak German quickly.}
            \end{itemize}
        \end{exbox}
    
    
    \section{Wichtige Verben mit Korrelat}
        
        \begin{grammarentry}{Typische Verben}{lernen als feste Struktur}
            \GermanText{Viele deutsche Verben sollte man zusammen mit ihrer Präposition und ihrem Korrelat lernen.}
            \EnglishText{Many German verbs should be learned together with their preposition and their correlative.}
            
            \begin{itemize}
                \item \textbf{sich freuen auf + Akk.} $\rightarrow$ \textbf{sich darauf freuen, ...}
                \item \textbf{denken an + Akk.} $\rightarrow$ \textbf{daran denken, ...}
                \item \textbf{warten auf + Akk.} $\rightarrow$ \textbf{darauf warten, ...}
                \item \textbf{achten auf + Akk.} $\rightarrow$ \textbf{darauf achten, ...}
                \item \textbf{rechnen mit + Dat.} $\rightarrow$ \textbf{damit rechnen, ...}
                \item \textbf{sorgen für + Akk.} $\rightarrow$ \textbf{dafür sorgen, ...}
                \item \textbf{sich beschäftigen mit + Dat.} $\rightarrow$ \textbf{sich damit beschäftigen, ...}
                \item \textbf{sich ärgern über + Akk.} $\rightarrow$ \textbf{sich darüber ärgern, ...}
            \end{itemize}
        \end{grammarentry}
        
        \begin{exbox}
            \begin{itemize}
                \ExampleItem{Ich warte darauf, dass du antwortest.}{I am waiting for you to answer.}
                \ExampleItem{Ich achte darauf, langsam zu sprechen.}{I pay attention to speaking slowly.}
                \ExampleItem{Ich beschäftige mich damit, bessere Sätze zu schreiben.}{I am dealing with writing better sentences.}
                \ExampleItem{Er ärgert sich darüber, dass der Bus wieder zu spät kommt.}{He is annoyed that the bus is late again.}
            \end{itemize}
        \end{exbox}
    
    
    \section{Wichtige Adjektive mit Korrelat}
        
        \begin{grammarentry}{Adjektiv + Präposition}{feste Verbindung}
            \GermanText{Auch viele Adjektive brauchen eine feste Präposition. Dann benutzt man oft ein da-Korrelat.}
            \EnglishText{Many adjectives also need a fixed preposition. Then German often uses a da-correlative.}
            
            Beispiele:
            \begin{itemize}
                \item \textbf{interessiert an + Dat.} $\rightarrow$ \textbf{daran interessiert}
                \item \textbf{stolz auf + Akk.} $\rightarrow$ \textbf{darauf stolz}
                \item \textbf{zufrieden mit + Dat.} $\rightarrow$ \textbf{damit zufrieden}
                \item \textbf{verantwortlich für + Akk.} $\rightarrow$ \textbf{dafür verantwortlich}
                \item \textbf{froh über + Akk.} $\rightarrow$ \textbf{darüber froh}
            \end{itemize}
        \end{grammarentry}
        
        \begin{exbox}
            \begin{itemize}
                \ExampleItem{Ich bin daran interessiert, Deutsch besser zu verstehen.}{I am interested in understanding German better.}
                \ExampleItem{Ich bin stolz darauf, dass ich jeden Tag übe.}{I am proud that I practice every day.}
                \ExampleItem{Ich bin damit zufrieden, wie viel ich gelernt habe.}{I am satisfied with how much I have learned.}
                \ExampleItem{Du bist dafür verantwortlich, pünktlich zu kommen.}{You are responsible for coming on time.}
            \end{itemize}
        \end{exbox}
    
    
    \section{Unterschied zwischen \textit{es} und da-Korrelaten}
        
        \begin{grammarentry}{es oder darauf / daran / damit?}{wichtiger Unterschied}
            \GermanText{\textbf{Es} benutzt man oft, wenn kein festes Präpositionsverb im Hauptsatz steht. Da-Korrelate benutzt man oft, wenn das Verb oder Adjektiv eine feste Präposition braucht.}
            \EnglishText{\textbf{Es} is often used when there is no fixed prepositional verb in the main clause. Da-correlatives are often used when the verb or adjective needs a fixed preposition.}
            
            Vergleich:
            \begin{itemize}
                \item \textbf{Ich finde es gut, Deutsch zu lernen.}
                \item \textbf{finden} braucht hier keine feste Präposition.
                \item \textbf{Ich freue mich darauf, Deutsch zu lernen.}
                \item \textbf{sich freuen auf} braucht die Präposition \textbf{auf}.
            \end{itemize}
        \end{grammarentry}
        
        \begin{exbox}
            \begin{itemize}
                \ExampleItem{Ich finde es wichtig, viel zu sprechen.}{I find it important to speak a lot.}
                \ExampleItem{Ich freue mich darauf, viel zu sprechen.}{I look forward to speaking a lot.}
                \ExampleItem{Ich halte es für sinnvoll, jeden Tag zu üben.}{I consider it useful to practice every day.}
                \ExampleItem{Ich denke daran, jeden Tag zu üben.}{I think about practicing every day.}
            \end{itemize}
        \end{exbox}
    
    
    \section{Position im Satz}
        
        \begin{grammarentry}{Wo steht das Korrelat?}{Wortstellung}
            \GermanText{Das Korrelat steht normalerweise im Hauptsatz. Der Infinitivsatz oder dass-Satz kommt danach.}
            \EnglishText{The correlative usually stands in the main clause. The infinitive clause or dass-clause comes after it.}
            
            Struktur:
            \begin{itemize}
                \item \textbf{Hauptsatz + Korrelat + Komma + Nebensatz}
                \item \textbf{Hauptsatz + Korrelat + Komma + Infinitivsatz}
            \end{itemize}
            
            Beispiele:
            \begin{itemize}
                \item \textbf{Ich finde es gut, dass du lernst.}
                \item \textbf{Ich finde es gut, Deutsch zu lernen.}
                \item \textbf{Ich freue mich darauf, dass du kommst.}
                \item \textbf{Ich freue mich darauf, dich zu sehen.}
            \end{itemize}
        \end{grammarentry}
    
    
    \section{Komma bei Korrelat-Sätzen}
        
        \begin{grammarentry}{Komma vor dass und zu}{sehr wichtig}
            \GermanText{Vor einem \textbf{dass-Satz} steht immer ein Komma. Vor einem erweiterten Infinitivsatz mit \textbf{zu} steht normalerweise auch ein Komma.}
            \EnglishText{Before a \textbf{dass-clause}, there is always a comma. Before an extended infinitive clause with \textbf{zu}, there is usually also a comma.}
            
            Beispiele:
            \begin{itemize}
                \item \textbf{Ich weiß, dass du kommst.}
                \item \textbf{Ich finde es gut, dass du kommst.}
                \item \textbf{Ich versuche, Deutsch zu sprechen.}
                \item \textbf{Ich finde es wichtig, jeden Tag Deutsch zu üben.}
            \end{itemize}
        \end{grammarentry}
    
    
    \section{Typische Fehler}
        
        \begin{grammarentry}{Fehler 1: falscher Artikel}{eine Frau}
            \GermanText{Bei femininen Nomen steht im Akkusativ \textbf{eine}, nicht \textbf{ein}.}
            \EnglishText{With feminine nouns, the accusative form is \textbf{eine}, not \textbf{ein}.}
            
            \begin{itemize}
                \item falsch: \textbf{Wenn ich ein Frau hätte}
                \item richtig: \textbf{Wenn ich eine Frau hätte}
            \end{itemize}
        \end{grammarentry}
        
        \begin{grammarentry}{Fehler 2: bevorzugen mit falscher Präposition}{kein in}
            \GermanText{Das Verb \textbf{bevorzugen} braucht normalerweise ein Akkusativobjekt ohne Präposition.}
            \EnglishText{The verb \textbf{bevorzugen} normally takes an accusative object without a preposition.}
            
            \begin{itemize}
                \item falsch: \textbf{Ich würde in ein Land bevorzugen.}
                \item richtig: \textbf{Ich würde ein Land bevorzugen.}
                \item besser: \textbf{Ich würde es bevorzugen, in einem Land zu leben.}
            \end{itemize}
        \end{grammarentry}
        
        \begin{grammarentry}{Fehler 3: Korrelat vergessen}{es / darauf / damit}
            \GermanText{Bei manchen Verben klingt der Satz ohne Korrelat unvollständig oder unnatürlich.}
            \EnglishText{With some verbs, the sentence sounds incomplete or unnatural without the correlative.}
            
            \begin{itemize}
                \item besser: \textbf{Ich freue mich darauf, dich zu sehen.}
                \item besser: \textbf{Ich rechne damit, dass er kommt.}
                \item besser: \textbf{Ich bevorzuge es, früh aufzustehen.}
            \end{itemize}
        \end{grammarentry}
        
        \begin{grammarentry}{Fehler 4: falsches Korrelat}{Präposition beachten}
            \GermanText{Das Korrelat hängt von der Präposition des Verbs oder Adjektivs ab.}
            \EnglishText{The correlative depends on the preposition of the verb or adjective.}
            
            Beispiele:
            \begin{itemize}
                \item \textbf{sich freuen auf} $\rightarrow$ \textbf{darauf}
                \item \textbf{denken an} $\rightarrow$ \textbf{daran}
                \item \textbf{rechnen mit} $\rightarrow$ \textbf{damit}
                \item \textbf{sorgen für} $\rightarrow$ \textbf{dafür}
            \end{itemize}
        \end{grammarentry}
    
    
    \section{Kurze Zusammenfassung}
        
        \begin{grammarentry}{Merksätze}{wichtig}
            \GermanText{Ein Korrelat ist ein kleines Wort, das auf einen späteren Satz verweist.}
            \EnglishText{A correlative is a small word that refers forward to a later clause.}
            
            \begin{itemize}
                \item \textbf{es} steht oft vor Infinitivsätzen und dass-Sätzen.
                \item \textbf{darauf, daran, damit, dafür, dazu} benutzt man oft bei Verben mit fester Präposition.
                \item Das Korrelat verweist oft nicht auf ein einzelnes Nomen, sondern auf eine ganze Handlung oder einen ganzen Satz.
                \item Bei Infinitivsätzen steht das Verb am Ende mit \textbf{zu}.
                \item Bei dass-Sätzen steht das konjugierte Verb am Ende.
                \item Die Präposition des Verbs bestimmt oft das richtige da-Korrelat.
            \end{itemize}
        \end{grammarentry}
        
        \begin{exbox}
            \begin{itemize}
                \ExampleItem{Ich bevorzuge es, in einem sicheren Land zu leben.}{I prefer to live in a safe country.}
                \ExampleItem{Ich freue mich darauf, Deutsch besser zu sprechen.}{I look forward to speaking German better.}
                \ExampleItem{Ich achte darauf, richtige Artikel zu benutzen.}{I pay attention to using correct articles.}
                \ExampleItem{Es ist wichtig, diese Strukturen oft zu üben.}{It is important to practice these structures often.}
            \end{itemize}
        \end{exbox}

\end{document}